\title{Large Language Models Can Automatically Engineer Features for Few-Shot Tabular Learning}

\begin{abstract}
Large Language Models (LLMs) have demonstrated extraordinary proficiency in handling complex and previously unencountered reasoning tasks, presenting significant opportunities for advancing tabular learning—a core requirement in many practical domains. In this study, we introduce \textsf{FeatLLM}, a novel in-context learning paradigm where LLMs are leveraged specifically for feature engineering, thus generating modified datasets optimized for tabular predictive modeling. These newly derived features serve as inputs to straightforward downstream machine learning algorithms, such as linear regression, enabling highly effective few-shot learning. Our \textsf{FeatLLM} method relies exclusively on this simplified predictive model paired with the LLM-generated features during inference, obviating the necessity for per-sample LLM invocations at test time that are commonplace in existing LLM-driven approaches. In addition, \textsf{FeatLLM} can operate with API-level access, sidestepping limitations on prompt size. Through extensive evaluation on various tabular datasets spanning multiple fields, \textsf{FeatLLM} not only yields high-quality feature rules but also achieves substantial performance gains (on average, a 10\% improvement) when compared to state-of-the-art baselines such as TabLLM and STUNT.
\end{abstract}

\section{Introduction}
Large Language Models (LLMs) have lately achieved notable success in generating relevant outputs for novel tasks without task-specific fine-tuning~\cite{creswell2022selection,imani2023mathprompter,taylor2022galactica}. These models, trained on massive textual resources, possess rich, generalizable prior knowledge that can often replace the necessity for hand-crafted expertise in diverse disciplines~\cite{Genomes2021,singhal2023large,wilcox2003role}. As such, they have facilitated automation across applications including conversational analysis~\cite{finch2023leveraging} and automated program repair~\cite{fan2023automated}. When provided with succinct task statements and a handful of representative examples, LLMs are capable of contextual understanding, allowing them to focus on salient features and key examples through prompt-based conditioning~\cite{brown2020language,zhao2023survey}. This underscores their ability to process data contextually, making LLMs highly suitable as advanced feature engineers.

There has been growing interest in transferring the reasoning and background knowledge encapsulated in LLMs to the problem of tabular learning. For example, prior domain insights, such as recognizing age as a risk factor for certain illnesses, can enhance tasks like disease outcome prediction. Some recent research has involved expressing task goals and feature meanings in natural language, followed by serializing input data and inputting it to an LLM~\cite{dinh2022lift,wang2023anypredict}. Subsequent steps typically incorporate either in-context learning~\cite{nam2023semi} or parameter-efficient finetuning approaches~\cite{dinh2022lift,hegselmann2023tabllm}. Empirical findings indicate that the use of LLMs’ embedded knowledge can consistently outperform traditional tabular modeling techniques, particularly in few-shot scenarios~\cite{hegselmann2023tabllm}.

Existing tabular learning approaches utilizing LLMs are constrained by several challenges. For end-to-end prediction scenarios, these methods necessitate at least a single LLM inference per data instance, substantially increasing computational demands. Achieving high predictive accuracy typically hinges on fine-tuning the LLM, which restricts the applicability to cases where the underlying LLM provides accessible training interfaces or tuning APIs. This presents a significant obstacle as state-of-the-art LLMs, highlighted in recent works~\cite{anil2023palm,openai2023gpt4}, often offer only restricted access through inference APIs. Moreover, such methods often fail to perform effectively with lengthy prompts~\cite{liu2023lost}; when the tabular data possesses a large number of features and the converted text input becomes extensive, these methods either suffer degraded performance or become impractical altogether. Collectively, these issues hinder the practical deployment of LLM-driven tabular learning in real-world contexts.

Our methodology diverges from conventional end-to-end prediction with LLMs by reassigning LLMs to the role of feature construction, thereby leveraging their pretrained knowledge more effectively. To this end, we propose a new in-context learning paradigm, \textsf{FeatLLM}, which seeks to uncover the implicit “criteria” used by LLMs in their prediction process. For example, in medical diagnosis, the LLM can deduce and articulate explicit rules connecting specific feature values to disease detection. These deduced criteria or rules are then harnessed to synthesize new features that serve as surrogates for the original features in downstream tasks. By repositioning the LLM as a feature generator rather than an inference engine, we enable the utilization of simpler downstream models—once features are generated, inference no longer requires complex models, substantially reducing latency. The in-context learning capabilities of LLMs allow for feature extraction via prompt-based example demonstrations, sidestepping any requirement for fine-tuning. Additionally, ensemble learning techniques such as feature bagging~\cite{breiman2001random} can be integrated to counteract issues arising from excessively long prompts.

\textsf{FeatLLM} decomposes the process of designing new features from prompts into two distinct stages: (1) interpreting the problem context and inferring how features relate to target variables, and (2) utilizing insights gained from the initial analysis, along with a limited number of labeled examples, to formulate precise rules that can separate different classes. This disciplined framework directs LLMs to ignore features lacking discriminative value when generating rules. These rules are subsequently enacted on the dataset (translating the rules into executable code), converting each instance into a binary attribute signifying compliance with every extracted rule. A simple predictive model is then fitted to use these derived binary features for class probability estimation. Robustness is further enhanced via an ensemble technique, which iteratively produces new features while leveraging feature bagging. Adjusting the quantity of features and instances sampled in each bag addresses the issue of overly lengthy prompts. The training-free nature of \textsf{FeatLLM} and its efficient inference make leveraging LLMs practical for a broad spectrum of real-world scenarios.

Our experimental assessment of \textsf{FeatLLM} covers 13 tabular datasets under low-data conditions, where it demonstrates consistently strong and stable performance. The results confirm that LLMs are adept at harnessing both embedded knowledge and empirical information from provided data, leading to the extraction of informative features. Across diverse tasks, our approach surpasses state-of-the-art few-shot learning methods. Implementation resources are made available through an anonymized GitHub repository at~\texttt{https://github.com/Sungwon-Han/FeatLLM}.

\section{Related Work}

\subsection{Few-Shot Learning with Tabular Data} 
Progress in few-shot learning methodologies has made it possible to obtain robust, transferable data representations while incurring only modest labeling expenses~\cite{chen2018closer}. Although the field began with a concentration on image datasets, its capacity for generalization has since been validated across several modalities such as text, audio, and most recently, tabular data~\cite{majumder2022few,nam2023stunt,schick2021exploiting}. A principal challenge in leveraging few-shot learning with limited labeled examples is the susceptibility to spurious associations~\cite{bartlett2021deep}. To mitigate this issue, contemporary strategies frequently integrate supplementary information or encode domain expertise as inductive biases to steer model development. These approaches include employing a surplus of accessible unlabeled data sets combined with thoughtfully designed augmentation to facilitate semi-supervised learning~\cite{ucar2021subtab,yoon2020vime}, or utilizing unsupervised pre-training methods that are agnostic to specific tasks to achieve advantageous model initializations~\cite{somepalli2021saint}. Unsupervised training objectives often consist of methods such as data masking or introducing corruption, subsequently using reconstruction losses~\cite{arik2021tabnet,majmundar2022met}, or adopting data augmentation frameworks for contrastive learning~\cite{bahri2022scarf,chen2020simple}. Nonetheless, the intrinsic heterogeneity present in tabular data complicates the identification of generic augmentations, and empirical evidence suggests these augmentation-based techniques fall short under few-shot regimes~\cite{nam2023stunt}.

More recently, approaches have advocated pre-training models on a diverse set of tabular datasets unrelated to the downstream tasks, subsequently applying transfer learning techniques to adapt to the target dataset~\cite{zhang2023generative,levin2022transfer,zhu2023xtab}. For example, TransTab~\cite{wang2022transtab} leverages transformer architectures to model semantic associations among table columns based on their names instead of solely relying on cell values, leading to transferability across a spectrum of tabular inference scenarios, including zero- and few-shot contexts. Conversely, TabPFN~\cite{hollmann2022tabpfn} synthesizes mixed-distribution data for the pre-training phase of a Bayesian neural network. At inference, the pre-trained model processes both training and test data from the target problem without further learning. The broad-spectrum knowledge drawn from multiple sources results in improved performance, outperforming conventional tree-based methods, particularly in data-sparse circumstances.

\subsection{Language-Interfaced Tabular Learning}
A further direction harnesses the prior knowledge embedded in large language models (LLMs)~\cite{anil2023palm,openai2023gpt4}. Due to their exposure to large, varied textual datasets, LLMs possess a remarkable aptitude for adapting to new tasks, making them broadly effective across many applications~\cite{brown2020language}. Recent research has endeavored to exploit this pre-trained knowledge base for tabular analysis by converting structured data into textual representations and embedding these as prompts for LLMs~\cite{dinh2022lift,hegselmann2023tabllm,wang2023anypredict}. Prompts enriched with table contents, feature annotations, and descriptions of the tasks are used, allowing LLMs to infer and reason about the underlying task-feature relationships necessary for prediction. This avenue has been explored through methods that either fine-tune model parameters in a resource-efficient manner, such as LoRA~\cite{hu2021lora} or IA3~\cite{liu2022few}, or capitalize on in-context learning by simply appending a handful of labeled examples to prompts, thereby sidestepping the need for model retraining~\cite{dinh2022lift,hegselmann2023tabllm,nam2023semi,slack2023tablet}.

Although previous techniques have demonstrated strong performance in few-shot tabular learning, their reliance on routing every inference query through the LLM limits their usability in real-world, industrial settings. This paper seeks to move beyond the traditional black-box pipeline, wherein predictions for each data point are produced via the LLM in a strictly end-to-end manner. Instead, the approach here employs the LLM to explicitly deduce decision conditions or rules for each target class. \textsf{FeatLLM} then translates these generated rules into executable program code, which, in turn, is used to engineer additional features. This permits predictions to be performed without further invocation of the LLM, leveraging the model’s capacity for in-context learning to recover rules from a combination of prior information and a limited number of labeled examples.

\section{Methods}
The core objective of \textsf{FeatLLM} is to identify ``rules''—that is, the defining criteria linked to each output class—from the handful of samples provided as input, as visualized in Figure~\ref{fig:main_model}. These rules are extracted via an LLM, subsequently converted into programmatic representations, and used to derive new binary features for each sample that indicate if specific rules hold. These binary features are then aggregated through a dot product with non-negative trainable weights to calculate the probability associated with each class. During training, the model internalizes the significance and effect of individual features on class assignment (see Section~\ref{Sec:training}). The rule extraction and feature creation processes are iterated in a bagging framework and integrated via ensemble techniques. Further clarifications on the problem setup and stepwise methodology are provided below.

\vspace*{-0.12in}
\paragraph{Problem formulation.} Consider a tabular dataset consisting of $N$ annotated instances, denoted as $\mathcal{D} = \{(\mathbf{x}^i, \mathbf{y}^i)\}_{i=1}^{N}$. Each data point $\mathbf{x}^i$ resides in a $d$-dimensional feature space, and every feature is associated with a natural language name from the set $F = \{f_j\}_{j=1}^{d}$; detailed definitions for these features are provided separately. In a classification scenario, the label $\mathbf{y}^i$ identifies the class to which the sample belongs. For $k$-shot learning setups, we randomly select $k$ ($<N$) labeled examples for model training.

\subsection{Prompt Design for Extracting Rules}
\label{Sec:prompt_design}
To facilitate effective rule extraction by the LLM via enhanced reasoning capabilities, we structure the prompt to mirror the logical steps a domain expert would follow when engaging with tabular data. The prompt is composed of three principal segments (see Fig.~\ref{fig:prompt_for_rule} and Appendix~\ref{sec:example_rule}):

\vspace*{-0.12in}
\paragraph{Basic information description.} This segment furnishes the necessary context for addressing the problem (highlighted in orange in Fig.~\ref{fig:prompt_for_rule}). It outlines the task (including class labels), supplies feature names and descriptions, and incorporates demonstration examples. The task prompt is articulated as a direct inquiry (e.g., ``Given this patient's myocardial infarction complications record, is there evidence of chronic heart failure? Yes or no?"; see Appendix~\ref{sec:task_desc} for additional examples). Each feature is accompanied by its data type and, optionally, its definition. For categorical features, representative category values may also be mentioned. To demonstrate input-output correspondences, a handful of training instances are serialized into textual format alongside their actual labels. This serialization adheres to conventions established in prior research~\cite{dinh2022lift,hegselmann2023tabllm}:

\begin{align}
&\text{Serialize}(\mathbf{x}^i,\mathbf{y}^i, F) = \nonumber \\ 
&``\ f_1\ \text{is}\ \mathbf{x}^i_1.\ ... \ f_d\  \text{is}\  \mathbf{x}^i_d.\ \ \text{Answer:}\  \mathbf{y}^i\ ”
\end{align}

\vspace*{-0.12in}
\paragraph{Reasoning instruction.} Our objective is to improve the LLM's reasoning capabilities by supplying explicit reasoning instructions (refer to the blue-highlighted text in Fig.~\ref{fig:prompt_for_rule}). These instructions open with a sentence modeled after the chain-of-thought paradigm~\cite{wei2022chain}. Following this, we organize the reasoning process for the LLM into two distinct stages. Initially, in the first stage, the LLM is prompted to independently assess and hypothesize the potential causal links or associations between features and the task objective, solely utilizing its intrinsic general knowledge or common sense—explicit example demonstrations are intentionally withheld at this point. This approach is intended to elicit the LLM's latent understanding relevant to the task. Subsequently, in the second stage, the LLM integrates the provided examples with its previous inferences from stage one, formulating explicit classification rules for each label. Adopting this staged strategy mitigates the risk of the model relying on coincidental patterns in irrelevant columns, thereby orienting attention towards features of substantive importance. In order to avoid issues such as underfitting or inflating the prompt due to either too few or too many extracted rules, we impose a fixed cap on the number of rules generated per class (specifically, 10)\footnote{A discussion of rule quantity is available in Section~\ref{sec:analysis}.}. Furthermore, while generating rules, the LLM is instructed to utilize the type specification for each feature from the foundational information section, which helps ensure that rules are appropriately aligned with feature types.

\vspace*{-0.12in}
\paragraph{Response instruction.} To streamline the extraction and practical application of derived rules, we provide explicit guidance regarding the organization of the LLM’s output (see the yellow-highlighted text in Fig.~\ref{fig:prompt_for_rule}). The prompt delineates clear requirements concerning both the formatting of outputs for each answer category (i.e., how responses should be formatted) and the expected structure for individual rules (i.e., rule formatting conventions). Through these specifications, the LLM can make informed choices on how to present rules according to the associated feature types. In cases where instructions are less prescriptive, the LLM may use logical connectives like AND or OR to combine multiple rules. An example output generated via these guidelines can be found in Appendix~\ref{sec:outcome-rule}.

\subsection{Inferring Class Likelihood via Rules}
\label{Sec:training}

\paragraph{Parsing rules for feature generation.} We employ the rules formulated in the earlier phase to derive new binary features. For every class, these features capture whether a given sample meets the criteria specified by the corresponding rules (represented as either '0' or '1'). Such features can subsequently inform class probability predictions. However, since the rules provided by the LLM are expressed in natural language, an effective parsing strategy is essential for automated feature extraction. While we standardize both response structure and rule presentation during rule generation to facilitate parsing, the LLM's outputs can still manifest unpredictable syntactic variations~\cite{kovavc2023large}. For example, the LLM might rephrase comparisons, using expressions like ``greater than'' or ``less than'' in place of ``$>$" or ``$<$", increasing parsing difficulty.

Rather than constructing elaborate programmatic parsers to handle such noisy variations, we instead exploit the LLM’s strengths. The LLM can often comprehend the intended semantics even when the syntax shifts and is capable of producing straightforward code when prompted (owing to its exposure to code-related data during training). To tap into these abilities, we provide prompts containing the function’s name, a clear description of inputs and outputs, along with the relevant rules. These prompts are then passed to the LLM for code generation. Figures~\ref{fig:prompt_for_parsing} and~\ref{sec:example_parsing}-\ref{sec:example_parse_outcome} in the Appendix illustrate sample prompts and their resulting outputs. The generated code is then executed in Python using the \texttt{exec()} function to facilitate the actual data transformation.

\vspace*{-0.12in}
\paragraph{Inferring class likelihood.} Upon obtaining binary features derived from the rules associated with each class, a straightforward approach for estimating the likelihood that a particular sample belongs to a class is to tally the number of rules from that class which are satisfied (that is, summing the binary features across all rules for each class). Yet, since individual rules and their corresponding features vary in significance, it becomes necessary to determine their respective importance based on information from the training data. To achieve this, we employ a conventional machine learning (ML) model to learn class-wise feature importance, applied individually to the binary features of each class. A linear model without a bias term can serve this purpose. Denoting the vector of binary features corresponding to sample $\mathbf{x}^i$ and class $k$ as $\mathbf{z}_k^i \in \{0, 1\}^R$—where $R$ stands for the number of class-specific rules and $c$ denotes the total number of classes—the class probability vector $\mathbf{p}^i$ is calculated by:

\begin{align}
\text{logit}_k^i &= \max(\mathbf{w}_k, 0) \cdot \mathbf{z}_k^i, \nonumber \\
\mathbf{p}^i &= \text{Softmax}({[\text{logit}_{1}^i, ..., \text{logit}_{c}^i]}). \label{eq:infer_prob}
\end{align}

Here, $\mathbf{w}_k$ comprises the learnable weights for the linear model, quantifying the importance assigned to each binary feature. By restricting the weights to the non-negative domain, we ensure that no rule discovered by the LLM has a negative impact on the estimated class likelihood during training. The logit for each class, derived this way, is subsequently converted to class probabilities through a softmax transformation. The parameters $\mathbf{w}_k$ are then learned by minimizing cross-entropy loss on a limited number of labeled instances (few-shot learning). 

\vspace*{-0.12in}
\paragraph{Ensembling with bagging.} In the final step, we repeat the overall procedure multiple times to train several inference models, whose outputs are aggregated to yield a consensus prediction. Specifically, given the ensemble of weight vectors trained over $T$ independent runs, denoted by $\{\mathbf{w}[t]\}_{t=1}^T$, the resulting prediction for a new sample is generated by averaging class probability vectors $\mathbf{p}^i$ as computed by Eq.~\ref{eq:infer_prob} for each $\mathbf{w}[t]$.

To incorporate randomness into rule generation within the ensemble, we set a relatively high temperature parameter during LLM inference. In addition, we shuffle the sequence of few-shot examples because previous findings indicate that demonstration order can affect LLM outputs~\cite{lu2022fantastically}\footnote{Appendix~\ref{sec:diversity} provides an analysis of rule diversity.}. To capitalize on predictive diversity and achieve more robust performance, we employ a bagging strategy, which samples subsets of features or instances for each iteration—a technique that becomes particularly valuable when dealing with large numbers of training examples or variables. Our ensemble methodology presents two key benefits. First, even when some trials produce suboptimal or incorrect rules, accurate rules generated by other ensemble members can offset these errors, collectively improving the resilience of the system. This advantage is supported by the principle of LLM self-consistency~\cite{wang2022self}, since rules consistently derived from multiple runs are more likely to be correct. Second, our ensemble method mitigates the challenge posed by LLM prompt size limitations. In cases where the volume of tabular data exceeds the model’s input length, bagging enables the reduction of input size while maintaining strong predictive results. Moreover, while a single rule set may not fully represent all underlying feature relations, aggregating multiple weak learners—each created with varied feature or instance coverage—ameliorates the limitations inherent in individual learners.

\section{Experiments}
We assess \textsf{FeatLLM} using several tabular datasets, conducting comprehensive analyses to better understand our framework's dynamics. Supplementary experiments, such as testing different LLMs or carrying out qualitative evaluations, are reported in the Appendix due to length restrictions.

\subsection{Performance Evaluation}
\paragraph{Datasets.}
We utilize 13 different datasets spanning binary and multi-class classification: (1) Adult~\cite{asuncion2007uci}, aiming to predict whether an individual earns more than \$50,000 per year; (2) Bank~\cite{moro2014data}, which involves forecasting whether a customer will open a term deposit; (3) Blood~\cite{yeh2009knowledge}, focusing on the likelihood of donor return; (4) Car~\cite{kadra2021well}, which assesses the quality of cars; (5) Communities~\cite{misc_communities_and_crime_183}, estimating crime rates for specific locales; (6) Credit-g~\cite{kadra2021well}, identifying if a person is a good or bad credit risk; (7) Diabetes\footnote{\texttt{kaggle.com/datasets/uciml/pima-indians-diabetes-database}}, predicting diabetes incidence; (8) Heart\footnote{\texttt{kaggle.com/datasets/fedesoriano/heart-failure-prediction}}, estimating the risk of coronary artery disease; and (9) Myocardial~\cite{misc_myocardial_infarction_complications_579}, predicting chronic heart failure outcomes.

Additionally, we incorporate several recently released and synthetic datasets that are infrequently examined in previous research and were excluded from LLMs’ pre-training corpora: (10) Cultivars, which requires predicting the potential grain yield of a specific soybean cultivar\footnote{\texttt{archive.ics.uci.edu/dataset/913}}; (11) NHANES, tasked with identifying the age group of individuals based on provided records\footnote{\texttt{archive.ics.uci.edu/dataset/887}}; (12) Sequence-type, where the objective is to classify a sequence as arithmetic, geometric, Fibonacci, or Collatz; and (13) Solution-mix, in which the aim is to determine whether the concentration of a four-solution mixture surpasses 0.5. The first two datasets were made publicly available in the UCI Machine Learning Repository after September 2023, guaranteeing they are absent from the training data of the LLM API (gpt-3.5-turbo-0613) leveraged by our model. The latter two are constructed artificially, preventing any prior exposure within the LLM API and thereby validating the integrity of subsequent evaluations.

A comprehensive summary of the datasets, covering their dimensions and degrees of complexity, is presented in Table~\ref{Tab:basic-desc}.
Each dataset features clearly labeled attributes accompanied by informative descriptions. For instance, datasets such as ‘Communities’ and ‘Myocardial’ contain over 100 attributes, introducing specific challenges given the constrained context window inherent to LLMs.

\vspace*{-0.12in}
\paragraph{Baselines.}
Our empirical assessment benchmarks against nine distinct baselines. The first set comprises established approaches in tabular learning: (1) Logistic regression (LogReg), (2) XGBoost~\cite{chen2016xgboost}, and (3) RandomForest~\cite{ho1995random}. The following two baselines exploit settings with only unlabeled data, specifically: (4) SCARF~\cite{bahri2022scarf}, and (5) STUNT~\cite{nam2023stunt}. It should be recognized that, in many realistic environments, amassing sufficient quantities of unlabeled data may not be feasible. The next comparator, (6) TabPFN~\cite{hollmann2022tabpfn}, constructs a synthetic distribution to pretrain models for improved tabular generalization. The final three baselines all utilize LLMs, namely (7) In-context learning (In-context)~\cite{wei2022emergent}, (8) TABLET~\cite{slack2023tablet}, and (9) TabLLM~\cite{hegselmann2023tabllm}. In-context learning involves incorporating a limited set of training examples directly into the input prompt, forgoing model adaptation. TABLET enhances the prompt via supplementary external cues—such as rule sets and prototype samples—provided by a separate classifier to boost reasoning accuracy. TabLLM uses parameter-efficient techniques, like IA3~\cite{liu2022few}, to adjust LLM parameters based on few-shot examples.

\vspace*{-0.12in}
\paragraph{Implementation details.}
\textsf{FeatLLM} employs GPT-3.5 as its LLM foundation, but the framework itself does not depend on any particular LLM (additional experiments using alternate backbones are presented in Appendix~\ref{sec:backbones}). For LLM inference, we specify a temperature value of 0.5 and retain the API’s default top-p setting of 1. The ensemble size and the count of extracted rules are set to 20 and 10, respectively. Further analysis concerning hyper-parameter choices can be found in Figure~\ref{fig:hyperparameter2} and Appendix~\ref{sec:hyper-parameter}.
For the linear model, we utilize the Adam optimizer with a learning rate of 0.01, training over 200 epochs. Selection of the optimal epoch is performed through $k$-fold cross-validation. When reproducing baselines, in-context learning utilizes GPT-3.5, while TabLLM follows its original configuration and employs T0~\cite{sanh2022multitask}. It should be noted that TabLLM restricts LLMs to those with openly released checkpoints, precluding the use of GPT-3.5. With regard to classical machine learning algorithms—such as LogReg, XGBoost, and RandomForest—hyperparameters are tuned using grid search in tandem with $k$-fold cross-validation. The value for $k$ is set to either 2 or 4 to ensure each class has representation in the training splits. Comprehensive implementation information is available in Appendix~\ref{sec:baselines}.

\vspace*{-0.12in}
\paragraph{Results.}
Tables~\ref{tab:main1} and \ref{tab:main2} display the performance metrics evaluated across multiple datasets. All experiments were independently repeated three times, each with a distinct split for training and testing sets, and the reported AUC values are represented as means with standard deviations. \textsf{FeatLLM} predominantly emerges as either the top or second-best method among the baselines considered. Remarkably, our approach matches the efficacy of standard tabular learning baselines using only 4-shot learning, compared to the latter's requirement for 32 shots (refer to Fig.~\ref{fig:compares}). While STUNT and TabPFN show notable advancements on several datasets, their overall gains relative to traditional machine learning algorithms are moderate. Furthermore, large-scale tabular datasets with many features present significant inference challenges for LLM-based approaches such as in-context learning, TABLET, and TabLLM, preventing them from generating valid results. In contrast, our method—which leverages both bagging and ensemble strategies—demonstrates robustness and effectiveness, even with high-dimensional feature spaces. Notably, integrating our bagging and ensemble methodology with alternative LLM frameworks is conceivable; however, this would effectively double the computation required per inference sample, resulting in prohibitive computational costs that undermine practical applicability. 

\subsection{Analysis \& Discussion}
\label{sec:analysis}
\paragraph{Ablation study.} We conduct ablation studies to isolate the contributions of each principal component to our system's performance, examining: (1) \textsf{-Tuning}: excluding weight tuning in the linear model and simply summing binary features to obtain the logit; (2) \textsf{-Ensemble}: removing the ensemble phase; (3) \textsf{-Description}: omitting the feature descriptions; and (4) \textsf{-Reasoning}: excluding the Step-1 reasoning instruction step during rule generation.

A summary of AUC changes from each ablation, averaged over all datasets, is provided in Table~\ref{tab:ablation}. Every modification to an individual component leads to reduced performance. Interestingly, the advantage of weight tuning becomes increasingly significant as the shot count rises, suggesting that an abundance of data allows for more precise rule weighting and thus enhances the impact of tuning. In contrast, the presence of feature descriptions and explicit reasoning instructions proves more influential under low-shot conditions, highlighting the value of effectively leveraging the inherent prior knowledge within the LLM. Finally, employing our ensemble methodology yields consistent and reliable improvements in predictive accuracy across all configurations.

\vspace*{-0.12in}
\paragraph{Training \& inference time comparison.} We analyze and contrast the training and inference durations for various methods. Experiments are conducted on the Adult dataset with a training set consisting of 16 samples. For all methods, one A100 GPU is used as standard, except for TabLLM, which operates over four A100 GPUs leveraging model parallelism. Detailed runtime metrics can be found in Table~\ref{tab:cost}. Importantly, \textsf{FeatLLM} achieves a notably low inference time, on par with established classical techniques. Conversely, inference with other LLM-driven methods, including In-context learning, TABLET, and TabLLM, incurs a substantially higher time cost. Regarding training time, \textsf{FeatLLM} matches other LLM-centric approaches, requiring merely 30 API calls—a quantity that does not significantly affect financial resources and which is amenable to parallel execution.

\vspace*{-0.12in}
\paragraph{Quality of features generated by \textsf{FeatLLM}.}
A straightforward experimental setup is implemented to examine how informative the rule-based features produced by \textsf{FeatLLM} are for real-world tasks. Concretely, the experiment calculates the average AUROC between each of the generated features and the ground truth labels, followed by determining the fraction of useful features—defined as those achieving AUROC above 0.5—across different datasets. The summary of these findings appears in Table~\ref{tab:quality}. The data clearly indicate that most generated rules have substantive association with the task-relevant labels and contribute meaningful information (see the ``Before tuning'' column). Further, throughout linear model optimization, any rules exhibiting minimal data relevance (i.e., those assigned weights falling below a predefined threshold) are excluded automatically (see the ``After tuning'' column). The threshold was fixed at 0.1, informed by the overall weight histogram.

\vspace*{-0.12in}
\paragraph{Effect of adding spurious correlations.} To evaluate how well different approaches mitigate the influence of extraneous spurious correlations in situations with limited labeled data, we carried out a targeted investigation. Our experiment employed the Heart dataset and randomly appended features from the Adult dataset—columns with no direct connection to the primary predictive task in Heart. By gradually increasing the number of these irrelevant columns, we systematically tracked model performance. To ensure comparability across methods, all models were provided with exactly 8 labeled examples and no unlabeled data, thus precluding the inclusion of techniques such as SCARF or STUNT in this analysis.

Figure~\ref{fig:spurious} presents how each method's performance responds to the presence of these artificial spurious correlations. Among all models tested, \textsf{FeatLLM} exhibited the smallest drop in accuracy, suggesting greater resilience to noisy or irrelevant features. This advantage likely arises from the framework’s use of prior knowledge to explicitly relate features to the task during rule extraction, thereby disallowing rules that depend on unrelated columns and promoting more reliable performance. Quantitatively, rules associated with these non-informative features constituted merely 1–2\% of the total set. Nonetheless, it is important to note that \textsf{FeatLLM} is also susceptible to learning spurious dependencies and misattributing causal relations when appended columns originate from domains more closely tied to the Heart task (refer to Appendix~\ref{sec:spurious-appendix} for detailed observations).

\vspace*{-0.12in}
\paragraph{Hyper-parameter analysis}
There are two principal hyper-parameters within \textsf{FeatLLM}: the ensemble count and the number of rules retrieved per class from each LLM query. To evaluate their influence, we ran a series of controlled experiments. The data revealed that as the ensemble number increases, performance shows an upward trend but eventually levels off, with improvement becoming marginal beyond a threshold (approximately 20), as demonstrated in Fig.~\ref{fig:appendix-hyperparameter1} (see Appendix). Regarding the rule count, our analysis found no statistically significant variation; however, selecting 10 rules offered an optimal compromise between the prompt’s length and overall performance (see Fig.~\ref{fig:hyperparameter2}). We hypothesize that increasing the number of rules raises input dimensionality, which, on one hand, enhances information but also introduces risks of overfitting, particularly in the low-shot learning context.

\section{Conclusion}
This work introduces \textsf{FeatLLM}, a method that sets a new benchmark for LLM-enabled few-shot learning on tabular data. Rather than relying solely on LLMs for direct prediction at the sample level, \textsf{FeatLLM} employs them as tools for deriving predictive criteria, akin to feature engineering. Integrating rule extraction with an ensemble-based bagging procedure, this framework offers low inference costs, is reliant only on API interaction (avoiding traditional model retraining), and circumvents data feature dimensionality constraints. These qualities allow \textsf{FeatLLM} to attain strong predictive outcomes, positioning it as an excellent choice for leveraging LLMs with tabular data in situations where data is scarce.

Although \textsf{FeatLLM} is currently restricted to low-shot learning environments, future work will target its adaptation for larger datasets by devising novel feature construction strategies that surpass mere rule formulation, along with enhancing interpretability. Such advancements would facilitate its broader application across a more diverse set of challenges.

\section{Broader Impact}
The framework described here, \textsf{FeatLLM}, is devised to leverage the expansive prior knowledge embedded in LLMs for tabular learning problems, offering performance benefits that traditional supervised approaches might not attain—especially in data-constrained environments. Our approach enhances data efficiency by tapping into this prior knowledge, thereby mitigating overfitting risks when labeled data is sparse. The method promises practical advantages for end users in areas like finance and healthcare, sectors where acquiring ground truth labels can be prohibitive due to labor or logistical barriers, and where the knowledge ingrained in LLMs can be particularly pertinent (since LLMs are typically pretrained on domain-relevant text corpora). Moving toward mainstream deployment, it is imperative to assess the implications of any social biases and potential misinformation transmitted via the LLMs’ prior knowledge, given the possibility that such embedded knowledge could substantially shape, or even supersede, model behavior in relation to observed data.

\end{document}